\begin{pastexam}[label={px:ntu-109-dsaa-15}]{臺大資工所 109 學年度 第 15 題}
（問答）

(15 points) A Single Nucleotide Polymorphism (SNP, pronounced \textit{snip}) is a single nucleotide
variation in the genome that recurs in a significant proportion of the population of a species.
The patterns of \textit{Linkage Disequilibrium} (LD) observed in the human population reveal a
block-like structure. LD refers to the association that particular alleles at nearby sites are
more likely to occur together than would be predicted by chance. The entire chromosome can be
partitioned into high LD regions interspersed by low LD regions. The high LD regions are usually
called ``haplotype blocks,'' and the low LD ones are referred to as ``recombination hotspots.''
Since there is little or no recombination within a haplotype block, these SNPs are highly
correlated. Consequently, a small subset of SNPs, called tag SNPs or haplotype tagging SNPs, is
sufficient to categorize the haplotype patterns of the block. It has been shown that we can
recast the tag SNP selection problem as \textbf{Problem W}, which is, ``Given a universal set
$\mathcal{U}=\{u_1,u_2,\ldots,u_n\}$ and a family $\mathcal{F}=\{F_1,F_2,\ldots,F_m\}$ of subsets
of $\mathcal{U}$, find a minimum-size subset $\widehat{\mathcal{C}}$ of $\mathcal{F}$, such that
every element of $\mathcal{U}$ belongs to at least one subset in $\widehat{\mathcal{C}}$.''
\begin{enumerate}[label=(\alph*)]
\item (5 points) Prove or disprove that a greedy approach always delivers an optimal solution for
\textbf{Problem W}.
\item (10 points) Prove or disprove that \textbf{Problem W} is NP-complete.
\end{enumerate}
\end{pastexam}

\begin{pastsoln}
答案：(a) 不成立，貪心不一定最佳；(b) 成立（指它的判定版本；最佳化版本是 NP 難）。

題目前面一大段生物背景只是包裝，\textbf{Problem W} 就是（不加權的）集合覆蓋問題（資結書第~29.4~節）。

\textbf{(a) 不成立，舉反例}。 「greedy approach」指最自然的貪心（資結書虛擬碼~29.3，每個集合成本都是
$1$）：每一步挑「蓋住最多還沒被蓋住的元素」的集合，直到全部蓋住。 反例：
\[
\mathcal{U} = \{1, 2, 3, 4, 5, 6\},\quad F_1 = \{1, 2, 3\},\quad F_2 = \{4, 5, 6\},\quad F_3 = \{1, 2, 4, 5\}.
\]
\begin{itemize}
\item 第一步：$F_1, F_2, F_3$ 能蓋住的新元素分別是 $3, 3, 4$ 個，挑 $F_3$（唯一最多，沒有平手的問題）。
還沒蓋住 $\{3, 6\}$。
\item 第二步：$F_1$、$F_2$ 各蓋住一個新元素（$3$、$6$），挑其中一個，例如 $F_1$。 剩 $\{6\}$。
\item 第三步：只有 $F_2$ 能蓋住 $6$，挑 $F_2$。
\end{itemize}
貪心得到 $3$ 個集合；但 $\{F_1, F_2\}$ 兩個就蓋住 $\mathcal{U}$，而一個集合最多 $4$ 個元素、蓋不住 $6$ 個，
所以最佳解是 $2$。 貪心不是最佳的。

換成別的貪心規則也一樣：由 (b)，Problem W 是 NP 難的，如果某個多項式時間的貪心演算法永遠給出最佳解，就
得到 $\mathrm{P} = \mathrm{NP}$。 貪心能保證的是近似：資結書定理~29.5，它的解最多是最佳解的 $H_d \le
\ln n + 1$ 倍（$d$ 是最大集合的大小）。

\textbf{(b) 成立}。 NP 完備是對判定問題說的，Problem W 是最佳化問題，先寫成判定版本（資結書第~28.1~節）：

\begin{center}
SET-COVER：給 $\mathcal{U}$、$\mathcal{F}$ 和整數 $k$，問是否存在 $\mathcal{C} \subseteq \mathcal{F}$，
$|\mathcal{C}| \le k$，蓋住 $\mathcal{U}$ 的每個元素？
\end{center}
證明 SET-COVER 是 NP 完備的；能在多項式時間解 Problem W，就能解 SET-COVER（解出最小的大小，和 $k$
比），所以 Problem W 是 NP 難的。 照資結書引理~28.6 的步驟：在 NP 裡，再從一個已知的 NP 完備問題歸約
過來。

\textbf{SET-COVER $\in$ NP}。 憑證是 $\mathcal{C}$ 本身（最多 $k \le m$ 個集合的編號）。 驗證：檢查
$|\mathcal{C}| \le k$，再把 $\mathcal{C}$ 裡每個集合的元素標記起來，看 $\mathcal{U}$ 的每個元素是否都被標到，
時間 $O\bigl(n + \sum_i |F_i|\bigr)$，是輸入長度的多項式。

\textbf{VERTEX-COVER $\le_{\mathrm{p}}$ SET-COVER}。 VERTEX-COVER（給圖 $G = (V, E)$ 和 $k$，問有沒有
$\le k$ 個頂點碰到每一條邊）是 NP 完備的（資結書定理~28.9）。 給一個 VERTEX-COVER 的輸入 $\langle G, k
\rangle$，構造
\[
\mathcal{U} = E, \qquad \mathcal{F} = \{\, F_v : v \in V \,\}, \quad F_v = \{\, e \in E : e \text{ 碰到 } v \,\},
\qquad k \text{ 不變}.
\]
構造只要掃一次相鄰串列，$O(|V| + |E|)$ 時間。 正確性：
\begin{itemize}
\item 若 $S \subseteq V$ 是大小 $\le k$ 的頂點覆蓋，則 $\{F_v : v \in S\}$ 有 $\le k$ 個集合；每條邊 $e$ 至少
有一個端點 $v \in S$，$e \in F_v$，所以蓋住了 $\mathcal{U} = E$。
\item 反過來，若 $\{F_v : v \in S\}$ 是大小 $\le k$ 的集合覆蓋，則每條邊 $e$ 在某個 $F_v$（$v \in S$）裡，
也就是 $e$ 碰到 $v$，所以 $S$ 是大小 $\le k$ 的頂點覆蓋。
\end{itemize}
所以 $\langle G, k\rangle$ 的答案是「是」若且唯若構造出來的 $\langle \mathcal{U}, \mathcal{F}, k\rangle$ 的
答案是「是」。 SET-COVER 在 NP 裡而且是 NP 難的，是 NP 完備的。 這個構造就是資結書第~29.4~節說的「頂點
覆蓋是集合覆蓋的特例」。

例如三角形 $G$：$V = \{x, y, z\}$，$E = \{xy, yz, zx\}$，$k = 2$。 構造出 $\mathcal{U} = \{xy, yz, zx\}$、
$F_x = \{xy, zx\}$、$F_y = \{xy, yz\}$、$F_z = \{yz, zx\}$；$\{x, y\}$ 是頂點覆蓋，$\{F_x, F_y\}$ 也正好
蓋住 $\mathcal{U}$。
\end{pastsoln}
